\documentclass[aspectratio=169, 9pt]{beamer}
\usetheme{metropolis}

\usepackage{xeCJK}
\setCJKmainfont{PingFang SC}
\setCJKsansfont{PingFang SC}
\setCJKmonofont{PingFang SC}

\usepackage{fontspec}
\setmonofont{Menlo}[Scale=0.82]
\usepackage{tikz}
\usetikzlibrary{arrows.meta, positioning, calc, shapes.geometric}
\usepackage{graphicx}
\usepackage{booktabs}
\usepackage{array}
\usepackage{listings}
\usepackage{xcolor}

\graphicspath{{images/}}

% ── Color palette ──────────────────────────────────────────────────
\definecolor{primary}{HTML}{6C5CE7}
\definecolor{accent}{HTML}{00B894}
\definecolor{warning}{HTML}{E17055}
\definecolor{dark}{HTML}{2D3436}
\definecolor{light}{HTML}{DFE6E9}
\definecolor{muted}{HTML}{636E72}
\definecolor{codebg}{HTML}{F7F7FA}
\definecolor{erlred}{HTML}{D63031}
\definecolor{erlblue}{HTML}{0984E3}

% ── Beamer theme ───────────────────────────────────────────────────
\setbeamercolor{normal text}{fg=dark,bg=white}
\setbeamercolor{frametitle}{fg=white,bg=primary}
\setbeamercolor{progress bar}{fg=accent,bg=light}
\setbeamercolor{alerted text}{fg=warning}
\setbeamercolor{block title}{fg=white,bg=primary}
\setbeamercolor{block body}{fg=dark,bg=primary!7}
\setbeamercolor{block title example}{fg=white,bg=accent}
\setbeamercolor{block body example}{fg=dark,bg=accent!7}
\setbeamertemplate{itemize item}{\color{accent}$\blacktriangleright$}
\setbeamertemplate{itemize subitem}{\color{primary}$\bullet$}
\setbeamertemplate{enumerate item}{\color{accent}\insertenumlabel.}
\metroset{progressbar=frametitle, numbering=fraction}

% ── Macros ─────────────────────────────────────────────────────────
\newcommand{\keypoint}[1]{\textbf{\color{primary}#1}}
\newcommand{\good}[1]{\textbf{\color{accent!80!black}#1}}
\newcommand{\warn}[1]{\textbf{\color{warning}#1}}
\newcommand{\code}[1]{\texttt{#1}}
\newcommand{\imgslide}[3]{%
  \begin{frame}{#1}
    \begin{columns}[T,onlytextwidth]
      \column{0.57\textwidth}
      \centering
      \includegraphics[width=\linewidth,height=0.72\textheight,keepaspectratio]{#2}
      \column{0.39\textwidth}
      #3
    \end{columns}
  \end{frame}
}

\lstdefinelanguage{Erlang2}{
  keywords={module,export,application,supervisor,gen_server,gen_statem,case,of,end,fun,receive,after,true,false},
  sensitive=true,
  morecomment=[l]{\%},
  morestring=[b]",
  morestring=[b]'
}
\lstset{
  language=Erlang2,
  basicstyle=\ttfamily\scriptsize,
  keywordstyle=\color{primary}\bfseries,
  commentstyle=\color{muted}\itshape,
  stringstyle=\color{accent!70!black},
  backgroundcolor=\color{codebg},
  frame=single,
  framesep=1.5mm,
  rulecolor=\color{light},
  showstringspaces=false,
  breaklines=true,
  tabsize=2
}

\title{Erlang observer\_cli 入门}
\subtitle{终端环境下观察 BEAM 节点、application、process 和 ETS}
\author{}
\date{}

\begin{document}

\begin{frame}[plain]
  \begin{tikzpicture}[remember picture, overlay]
    \fill[primary!8] (current page.north west) rectangle (current page.south east);
    \fill[primary!15] (current page.north west) ++(1.7,-1.3) circle (2.7cm);
    \fill[accent!12]  (current page.south east) ++(-2.1,1.6) circle (3.0cm);
    \node[anchor=center, font=\fontsize{25}{31}\selectfont\bfseries, text=dark]
      at ([yshift=1.7cm]current page.center) {observer\_cli};
    \node[anchor=center, font=\Large, text=primary!80]
      at ([yshift=0.55cm]current page.center) {没有 GUI 时怎样观察 Erlang 节点};
    \node[anchor=center, font=\large\bfseries, text=accent!80!black]
      at ([yshift=-0.35cm]current page.center) {SSH、容器、服务器环境中的 BEAM 运行时视角};
    \draw[accent, line width=1.5pt]
      ([yshift=-1.15cm, xshift=-4.5cm]current page.center)
      -- ([yshift=-1.15cm, xshift=4.5cm]current page.center);
    \node[anchor=south, font=\scriptsize, text=dark!45]
      at ([yshift=0.4cm]current page.south)
      {基于 full\_otp\_app demo 与 observer\_tutorial\_body.tex 整理};
  \end{tikzpicture}
\end{frame}

\begin{frame}{Observer UI 和 observer\_cli 的关系}
  \begin{columns}[T,onlytextwidth]
    \column{0.48\textwidth}
    \begin{block}{Observer UI}
      \begin{itemize}
        \item 图形界面
        \item 适合本地开发和教学截图
        \item 强项是 application / supervisor tree
        \item 可以双击进入进程详情
      \end{itemize}
    \end{block}
    \column{0.48\textwidth}
    \begin{exampleblock}{observer\_cli}
      \begin{itemize}
        \item 终端界面
        \item 适合 SSH、容器、无桌面环境
        \item 强项是指标和排序
        \item 可以查看 process / ETS / system / network
      \end{itemize}
    \end{exampleblock}
  \end{columns}
  \vspace{0.3cm}
  \centering
  \good{两者都在观察同一个 BEAM 运行时，只是交互方式和重点不同。}
\end{frame}

\begin{frame}[fragile]{依赖和启动方式}
  \begin{columns}[T,onlytextwidth]
    \column{0.48\textwidth}
    \begin{block}{rebar.config}
\begin{lstlisting}
{deps, [
    {observer_cli, "1.8.8"},
    {recon, "2.5.6"}
]}.
\end{lstlisting}
    \end{block}
    \column{0.48\textwidth}
    \begin{block}{启动 demo 和 observer\_cli}
\begin{lstlisting}[language=bash]
cd bilibili_erlang/observer/full_otp_app
rebar3 shell --sname application_demo
\end{lstlisting}
\begin{lstlisting}
observer_cli:start().
\end{lstlisting}
    \end{block}
  \end{columns}
\end{frame}

\imgslide{Home：启动后的主界面}{observer_cli_start.png}{%
  \begin{block}{顶部菜单}
    \begin{itemize}
      \item \code{Home(H)}
      \item \code{Network(N)}
      \item \code{System(S)}
      \item \code{Ets(E)}
      \item \code{App(A)}
      \item \code{Doc(D)} / \code{Plugin(P)}
    \end{itemize}
  \end{block}
  \begin{exampleblock}{常用按键}
    \code{H} 回 Home，\code{A} 看 application，\code{E} 看 ETS，\code{q} 退出。
  \end{exampleblock}
}

\begin{frame}{Home 页面适合看什么？}
  \begin{columns}[T,onlytextwidth]
    \column{0.50\textwidth}
    \begin{block}{节点级信息}
      \begin{itemize}
        \item Process count
        \item Memory
        \item Scheduler / run queue
        \item Ports / ETS / atoms
      \end{itemize}
    \end{block}
    \column{0.46\textwidth}
    \begin{exampleblock}{进程列表}
      \begin{itemize}
        \item 按 memory 排序
        \item 按 reductions 排序
        \item 按 message queue 排序
        \item 输入行号进入进程详情
      \end{itemize}
    \end{exampleblock}
  \end{columns}
\end{frame}

\imgslide{App(A)：查看所有 application}{observer_cli_application.png}{%
  \begin{block}{App 页面展示聚合指标}
    \begin{itemize}
      \item \code{App}
      \item \code{ProcessCount}
      \item \code{Memory}
      \item \code{Reductions}
      \item \code{MsgQ}
      \item \code{Status} / \code{Version}
    \end{itemize}
  \end{block}
  \begin{alertblock}{适合看资源消耗}
    例如观察 \code{full\_otp\_app} 是否 Started，以及它占用多少进程和内存。
  \end{alertblock}
}

\begin{frame}{App 页面不能像 GUI 一样展开 application}
  \begin{columns}[T,onlytextwidth]
    \column{0.48\textwidth}
    \begin{block}{Observer GUI Applications}
      \centering
      \code{application}\par
      $\downarrow$\par
      \code{supervisor}\par
      $\downarrow$\par
      \code{worker process}\par
      \vspace{0.2cm}
      \good{适合讲结构}
    \end{block}
    \column{0.48\textwidth}
    \begin{alertblock}{observer\_cli App(A)}
      \centering
      \code{application}\par
      $\downarrow$\par
      \code{ProcessCount / Memory}\par
      \code{Reductions / MsgQ}\par
      \vspace{0.2cm}
      \warn{不能展开 supervisor tree}
    \end{alertblock}
  \end{columns}
  \vspace{0.25cm}
  \centering
  App 页面里的行号是表格编号，不是进入 application 的菜单项。
\end{frame}

\begin{frame}[fragile]{命令行查看监督树仍然用 Erlang API}
  \begin{block}{如果想在终端看 full\_otp\_app 的子进程}
\begin{lstlisting}
supervisor:which_children(full_otp_app_sup).
supervisor:which_children(data_store_sup).
\end{lstlisting}
  \end{block}
  \begin{exampleblock}{解释}
    \code{observer\_cli} 的 App 页面做 application 级别聚合。监督树结构仍然可以用 Erlang shell 或脚本直接调用 OTP API 查看。
  \end{exampleblock}
\end{frame}

\imgslide{Processes：进程列表和进程详情}{observer_cli_processes.png}{%
  \begin{block}{进程列表可以进入详情}
    \begin{itemize}
      \item \code{m} 按 memory 排序
      \item \code{r} 按 reductions 排序
      \item \code{mq} 按 message queue 排序
      \item 输入进程行号打开详情
    \end{itemize}
  \end{block}
  \begin{alertblock}{和 App 页面不同}
    Process 行号可以进入详情；App 行号不能进入 application 详情。
  \end{alertblock}
}

\begin{frame}{教学时重点找这些 demo 进程}
  \begin{columns}[T,onlytextwidth]
    \column{0.48\textwidth}
    \begin{block}{supervisor}
      \begin{itemize}
        \item \code{full\_otp\_app\_sup}
        \item \code{data\_store\_sup}
      \end{itemize}
    \end{block}
    \begin{block}{gen\_server}
      \begin{itemize}
        \item \code{kv\_server}
        \item \code{data\_store}
        \item \code{data\_cache}
      \end{itemize}
    \end{block}
    \column{0.48\textwidth}
    \begin{exampleblock}{其他 behaviour}
      \begin{itemize}
        \item \code{event\_bus}
        \item \code{traffic\_light}
      \end{itemize}
    \end{exampleblock}
    \begin{alertblock}{进程详情信息}
      registered name、links、monitors、message queue、dictionary、current stack、state。
    \end{alertblock}
  \end{columns}
\end{frame}

\imgslide{Ets(E)：查看 ETS 表指标}{observer_ets.png}{%
  \begin{block}{ETS 页面适合线上排查}
    \begin{itemize}
      \item 表数量
      \item 表大小
      \item owner process
      \item memory / object count
    \end{itemize}
  \end{block}
  \begin{exampleblock}{和 GUI Table Viewer 对比}
    GUI 适合教学看表数据，CLI 适合服务器终端看表指标。
  \end{exampleblock}
}

\begin{frame}{observer\_cli 背后不是魔法}
  \begin{block}{它本质上依赖 BEAM 可观测 API}
    \begin{itemize}
      \item \code{erlang:system\_info/1}
      \item \code{erlang:memory/0}
      \item \code{application:which\_applications/0}
      \item \code{erlang:processes/0}
      \item \code{erlang:process\_info/2}
      \item \code{ets:all/0} / \code{ets:info/1}
    \end{itemize}
  \end{block}
  \begin{exampleblock}{demo 中的辅助脚本}
    \code{full\_otp\_app/scripts/observer\_cli\_snapshot.escript} 直接调用这些 API 采集一次快照。
  \end{exampleblock}
\end{frame}

\begin{frame}[fragile]{远程节点参数：node name，不是 OS PID}
  \begin{block}{获取当前节点名和 cookie}
\begin{lstlisting}
node().
erlang:get_cookie().
\end{lstlisting}
  \end{block}
  \begin{block}{脚本远程采集形式}
\begin{lstlisting}[language=bash]
scripts/observer_cli_snapshot.escript application_demo@HOST COOKIE
\end{lstlisting}
  \end{block}
  \begin{alertblock}{关键点}
    参数是 Erlang node name 和 cookie，不是 BEAM 的操作系统进程 PID。
  \end{alertblock}
\end{frame}

\begin{frame}{什么时候用 observer\_cli？}
  \begin{columns}[T,onlytextwidth]
    \column{0.48\textwidth}
    \begin{block}{推荐场景}
      \begin{itemize}
        \item SSH 到服务器
        \item 容器内排查
        \item 没有图形界面
        \item 想快速按指标排序
        \item 看 ETS / process / memory
      \end{itemize}
    \end{block}
    \column{0.48\textwidth}
    \begin{alertblock}{不适合场景}
      \begin{itemize}
        \item 教初学者看 supervisor tree
        \item 希望像 GUI 一样展开 application
        \item 需要大量截图标注结构
      \end{itemize}
    \end{alertblock}
  \end{columns}
\end{frame}

\begin{frame}{课堂演示顺序}
  \small
  \begin{enumerate}
    \item 启动 \code{rebar3 shell --sname application\_demo}
    \item 执行 \code{observer\_cli:start().}
    \item 看 Home 页面，说明这是终端版运行时观察工具
    \item 按 \code{A} 看所有 application 的聚合指标
    \item 说明 App 页面不能展开 supervisor tree
    \item 回到 Home，按 memory / reductions / message queue 排序进程
    \item 输入进程行号查看进程详情
    \item 按 \code{E} 查看 ETS 页面
    \item 最后引出 Erlang API 和 \code{observer\_cli\_snapshot.escript}
  \end{enumerate}
\end{frame}

\begin{frame}{小结}
  \begin{block}{本集结论}
    \begin{itemize}
      \item \code{observer\_cli} 是终端环境下的 BEAM 观察工具
      \item App 页面看 application 聚合指标，不展开监督树
      \item Process 列表可以进入进程详情
      \item ETS 页面适合服务器排查表指标
      \item 它背后依赖 Erlang runtime introspection API
    \end{itemize}
  \end{block}
  \begin{exampleblock}{和 Observer UI 配合}
    UI 版讲结构，CLI 版讲线上观察和指标排查，两者互补。
  \end{exampleblock}
\end{frame}

\end{document}
